\section{The Emperor’s Post}

Iamblichus, Hercules, the \emph{Meditations of the Tarot}, and Mouravieff's system of numbers can be correlated, with study and practice. I am stopping (for now) at Number Four, the Emperor, and drawing more closely several threads together to illustrate how this is so.

Kukai\footnote{\url{http://plato.stanford.edu/entries/kukai/}} I have mentioned before as a co-inhabitor of tradition.

In the section 3.12 on developments of the state of mind, we find this interesting passage:

\begin{quotationx}
1\textsuperscript{st} to 3\textsuperscript{rd} states: Pre-Buddhist stages: worldly ``vehicles'' of \emph{samsaric} entrapment:

1\textsuperscript{st} state: ``The mind of the goat foolishly transmigrating in the six destinies (or realms)'' (\emph{ishô teiyô-shin}): The state of desire driven by animal instincts without moral restraint; the stage to which belong common people, hell-beings, hungry-ghosts, beasts, \emph{asuras} (``titans''), and various deities or celestial beings trapped in their \emph{samsaric} destinies.

2\textsuperscript{nd} state: ``The mind of the child tempered but ignorantly obsessed with moral precepts'' (\emph{gudô jisai-shin}): The state of ethical actions and virtue that promote social order but without any ``religious'' goal; the stage to which belong Confucianism and the Buddhist precepts (\emph{ritsu}) for the laity.

3\textsuperscript{rd} state: ``The mind of the child composed and fearing nothing'' (\emph{yodô mui-shin}): The state of deity worship and extrinsic magico-religious practice for the sake of overcoming anxiety with the thought of attaining supernatural powers or immortality, or reaching an eternal and blissful heaven; the stage to which belong Taoism and various forms of Hinduism or Brahmanism.

4\textsuperscript{th} to 10\textsuperscript{th} states: Buddhist stages (the fourth to ninth being exoteric Buddhism and the tenth being esoteric Buddhism):

4\textsuperscript{th} to 5\textsuperscript{th} states: Hinâyâna stages: ``vehicles'' of those who aspire towards self-enlightenment without caring for the enlightenment of others.

4\textsuperscript{th} state: ``The mind of one affirming only the elements and negating the self'' (\emph{yuiun muga-shin}): The state of the \emph{?râvaka} who analyzes phenomena into the psycho-physical ``aggregates'' (\emph{skandhas}) and/or the elements (\emph{dharmas}), to thus negate any belief in a permanent ego (\emph{atman}); the stage to which belong the teachings of the historical Buddha and his direct disciples and of the \emph{Abhidharma} scholastics. While the substantiality of reality is thus deconstructed into its elemental \emph{dharmas}, the \emph{dharmas} themselves however become fetters, thus taking from three lives to sixty aeons to achieve liberation.

5\textsuperscript{th} state: ``The mind freed from karmic seeds'' (\emph{batsu gôinju-shin}): The self-other duality and recognizing the interdependency between self-enlightenment and other-enlightenment and between wisdom and compassion.
\end{quotationx}


The reader of Mouravieff's \emph{Gnosis} will no doubt recognize that Man 1, Man 2, and Man 3 (of which man \#3, the man of rationalized intellect is dominant in our era), are excluded from enlightenment, just as the esoteric Orthodox tradition teaches through Mouravieff's work. That is, the man of movement, the man of heart, and the man of mind are all subject to samsaric entrapment, by virtue of the very virtue which characterizes them, the organizing of their being around a privileged center. It is Man \# 4, the man who has begun to balance his centers through the difficult process of self-knowledge (which is painful as well as arduous), who can be said to begin to ``inherit the kingdom of God'' (the basilieia of God\footnote{\url{https://en.wikipedia.org/wiki/Kingdom\_of\_God\_\%28Christianity\%29}}).

In a healthy civilization, the men of intellect would have supervised the building of the edifice of the Ekklesia, the Oikumene, or the Third Kingdom. This edifice would have centered around the natural revelation of a worldly order, a worldly order ultimately ruled through the self-sacrifice of the Emperor (who, significantly, is the Fourth Card of the Tarot, thus superseding or sitting in guidance over the first three, in relation to manifesting the King of the World). In this way, authority is allowed to ``come out into the open\footnote{\url{https://bonald.wordpress.com/2015/01/07/not-ready-for-monarchy/}}``, and to play its proper role in saving cosmic energy by ``conserving'' energy, thus freeing up an excess for potential use for divine purposes that go beyond maintaining the order of things. This is the mystery behind Romans 13 and other verses, in which the apostles admonished Christians ``to love the brotherhood, fear God, and honor the King'', and to pray for those in authority, who bear not the sword in vain, that they may live quiet and peaceable lives of godliness.

In the Fourth Labor of Hercules\footnote{\url{https://gornahoor.net/?p=6806}}, the hero inherited the power of the man-beast Chiron, who is unable to bear immortality due to the dreadful wound he received from Hercules. It is important to note Tomberg's astonishing claim that the Emperor is the most ``human of men'', and that this entitles him to assume the post of guardian of the shields of the earth. King David is the paragon of this fact -- herder, warrior, champion, outlaw, king, sinner, priest, and initiate: his life was a full one, and his heart for his son Absalom made him a man ``after God's own heart'', being the paternal instinct writ so large that it conquered even the hatred of his own seed. Indeed, we find in history that men like King Alfred, King Arthur, and all the other paragons of chivalry were ``men in full'' : not the one first to boast, or the flashiest, or the loudest, but the one last to abandon the cause of earth. It is this love of simplicity which gives the emperor his wounds and makes him the ``best of men'', worthy to create an age in which deeds even greater than his can be done. The card of the Emperor teaches the divine conservation of what it is to be human --- or, phrased exoterically, ``whatsoever a man sows, that shall he reap'' \& ``I have reserved 7000 who have not bowed the knee to Baal''. Thus, Man not only endures, he \emph{prevails}.

In all of these studies, a kind of ``evolutionary test'' (Mouravieff's terminology) is being meted out : the rule is, use what you have for the best, or forfeit those talents to someone else, for man is neither beast nor machine. If ancient Greece's common cultural task was to delineate Man from wine and centaurs and the wild boar, then our time's task is to insist that Man is separate from the machine.

To continue with Mouravieff, Man \#2 is easily lead astray by ideals, choosing ones which are either flawed, or suit his predilections and accidental being. Man \#1 is swiftly lead astray in a dark wood, wandering, of sensual attractions. Man \#3, when corrupted, spearheads the ideological annihilation of natural and traditional order. When not corrupted, the order that they organize intellectually and spiritually provide the possibility for all three orders to co-exist peacefully and creatively and to lasting profit. But in actual practice, Man\#4 represents any of the other three basic types which have begun to awake, and thus it is also true that a second Triad is generated to begin organizing in all serious earnest in a way that defeats the natural tendency of entropy. Thus, Guenon emphasizes, in his discussion on Time and the Circle in Chapter 23 of the Great Triad\footnote{\url{https://www.gornahoor.net/library/Rene\_Guenon\_The\_Great\_Triad.pdf}}, the role of the Cross, or the four quarters of the world. These quarterings are made manifest in the four wounds of the emperor, who may be a man of action, a man of heart, or a man of mind, provided he has followed the divine in achieving, through Love, a participation in each. The Three Orders then, fully function and are protected, under this Lord.

Distinctively, he achieves this in full immersion in secular affairs, as the ``most human of men'', and is thus not the greatest of Men \#4, but the one among them who is the most human, and therefore most fit to rule. He is a man with the ``gift of tears'', for instance, and one finds (correspondingly) a higher development of the idea of chivalry in the West than the East (although it is not absent, identity with distinctive difference).

If the Emperor sits on the Imperial Throne, there can be many Kings made straight, who are not fully at the level of Man \#4, but whose type resonate strongly enough for the energy of the Emperor to magnetize towards this eventual goal, or at least constrain within an outer semblance of that path.

It is hard to predict how this web works, because we often stand outside, but in practice, any Man \#4 who stays true enough will cause repercussive changes beneath him that tend to drag more of prime materia and Creation towards the source of his Love, whereas those beneath him will tend to deteriorate without his immediate secular influence and power, even if occulted.

Whereas the tendency of the Emperor is to accrue forfeited spiritual power, rather than rely upon external influences. This does not preclude the possibility of a corruption and fall, only to say the impulse is now on the side of Heaven, with the passage to Man \#4. It is a mistake to think that ``divine right of Kings'' taught the infallibility of annointed rulers, and use this as a straw man.

The ravaging Erymanthian Boar (``the wild hog in the woods\footnote{\url{https://www.gornahoor.net/library/Rene\_Guenon\_The\_Great\_Triad.pdf}}'' in the old folk tune passed down from the times of medieval England) is the ravage of Time, the entropy of Creation which is set loose from divine balance. It is in this way that the pope claimed that ``a wild boar is loose in the vineyard of God'' when it became evident that Luther was a man of character who was no paper tiger. It is significant that it was the princes of Germany and the king of England whose betrayal and subversion lead to a complete rupture between the Protestant and Catholic worlds.

Mouravieff explains that some of this came to pass because the character of the people is an element that has to be constantly adjusted to in government : whoever has aspirations to progress along the path laid out by the Emperor will have to manage the first Triad of forces, and become used to discerning the neutralizing force which mediates between active/passive. If this is done under the will of God, such an ability will tend to draw that person into the character of real \emph{Auctoritas}, in whatever circle they inhabit, big or small.


I will note, for now, the correspondence of the ``oscillations'' and the juggling motif of the first Tarot card with the center of movement, and the correspondence of the second\footnote{\url{https://gornahoor.net/?p=7287}} and third\footnote{\url{http://www.fourhares.com/pdfs/mott/Letter-3.pdf}} cards with the emotional and intellectual center, respectively. Much more can and should be said about these correspondences: for one thing, it is the achievement of an intellectual connection and balance which allows the conception of the God-man, who is (of course) the new King of Kings and Lord of Lords, and who opposes the purely (in relation to himself) material world of the Prince of this World.

But the emperor, he is at a post, in balance: a certain level of civilization, real civilization and not the fake stuff masquerading in rhetoric, has been achieved. The shields of the earth can be marshaled and drawn to their rightful destinies. He is not the most sensual, the most emotional, or the most intelligent man, but the most \emph{human, and therefore in a more complete sense, remembers what it means to enjoy the material world, yearn for the heart of all meaning, or see into the pattern of the Logos}. In a world that hates the cosmic Law, men find it easy to forget the totality of Man -- that is, until the emperor enters from stage right, when time stops again.


\flrightit{Posted on 2015-01-24 by Logres}

